%=======================================================================
%  CHAPTER 5 — EQUALIZATION AND DIVERSITY TECHNIQUES  (4 hrs)
%=======================================================================
\section{Equalization and Diversity Techniques}

{\color{sub}\footnotesize 4 hrs \gap$\cdot$\gap asked in \mk{21 of 22 papers}
\gap$\cdot$\gap 7 syllabus sub-topics \gap$\cdot$\gap 1 numerical family
$\rightarrow$ Ch 5 numerical companion.}

\lead{Read this first:}
\begin{itemize}
  \item \textbf{Equalization} \tS{15} and \textbf{diversity} \tS{11} are two of the three most
        reliable questions in the whole paper. \textbf{RAKE} \tS{6} is close behind.
  \item Open every answer here with the \textbf{same framing sentence}: equalization, diversity
        and channel coding are three techniques used \textbf{independently or in tandem} to
        improve receiver signal quality.
  \begin{itemize}
    \item \textbf{Equalization} undoes \textbf{ISI} from a time-dispersive channel.
    \item \textbf{Diversity} undoes \textbf{fading} by giving the receiver more than one look at
          the signal.
    \item \textbf{Channel coding} (Ch 6) adds \textbf{redundancy} so faded bits can be recovered.
  \end{itemize}
  \item \mk{Learn which impairment each one fixes.} Papers repeatedly ask ``why is X needed'',
        and the mark is for naming the impairment, not describing the block diagram.
  \item \textbf{No numerical question on this chapter has ever appeared in a PYQ}, but the MSE
        optimum-weight derivation \tF{3} is examined as bookwork, and the tutorial set carries
        diversity SNR problems. Both are in the companion.
\end{itemize}

\hr

\T{5.1 Equalization: why, and how it adapts}

\Q{Why is equalization needed? Describe the signal processing operation that minimizes ISI}{\m{2+6}\m{2+5}\m{4}\m{2+1+4}}
{\color{sub}\footnotesize \tS{6} \yr{72 Ma, 73 Ma, \textbf{75 Bh}, \textbf{73 Bh},
\textbf{\texttt{82 Bh}}, \texttt{80 Ba}} \gap$\cdot$\gap \mk{the highest-frequency question in
the chapter}}

\sbsr{0.53}{%
\lead{Why it is needed}
\begin{itemize}
  \item A mobile channel is \textbf{time dispersive}: multipath copies arrive at different delays.
  \item A delayed copy of one symbol lands \textbf{on top of the next} symbol
        $\rightarrow$ \textbf{intersymbol interference (ISI)}.
  \item ISI happens whenever signal bandwidth $W$ exceeds the coherence bandwidth $B_C$, ie
        whenever the channel is \textbf{frequency selective}.
  \item ISI is \mk{the single largest bottleneck on high-rate data over a mobile channel}, and
        raising transmit power does not help: it raises the interference too.
\end{itemize}}{%
\figT{c5_eq_block.png}
\penalty0\vspace{1pt}
{\color{sub}\footnotesize A communication system w/ an adaptive equalizer at the receiver.
$f(t)$ is the \textbf{combined} impulse response of transmitter, radio channel and receiver
RF/IF stages. The equalizer inverts all three at once.}}

\sbs{%
\lead{Definition}
\begin{itemize}
  \item \textbf{Equalization} $=$ \textbf{any signal processing operation that minimizes ISI}.
  \item The equalizer is the \textbf{inverse filter of the channel}. If the channel response is
        $f(t)$, the equalizer response $h_{eq}(t)$ satisfies
        \[ f(t) * h_{eq}(t) = \delta(t) \]
  \item So it \textbf{flattens} the frequency response, hence the name.
  \item Implemented at \textbf{baseband or at IF} in the receiver, where the complex envelope is
        available.
  \item \mk{The mobile channel is random and time varying}, so a fixed filter cannot work. The
        equalizer must track the channel $\rightarrow$ it must be \textbf{adaptive}.
\end{itemize}}{%
\lead{The error-driven loop}
\begin{itemize}
  \item Equalizer output $\hat{d}(t)$ is compared against a desired signal $d(t)$.
  \item Difference $=$ \textbf{error signal} $e(t)$, fed back to retune the filter.
  \item If $d(t)$ is \textbf{not} fed back to adapt the filter, the equalizer is \textbf{linear}.
  \item If $d(t)$ \textbf{is} fed back to change subsequent outputs, it is \textbf{nonlinear}.
        \mk{That single sentence is the definition of the linear / nonlinear split.}
\end{itemize}}

\Q{Explain the training and tracking modes of an adaptive equalizer}{\m{4}\m{1+3}}
{\color{sub}\footnotesize \tF{3} \yr{\textbf{78 Ch}, \textbf{73 Bh}, \texttt{81 Ba}}}

\sbs{%
\lead{Training mode \mk{(measure the channel)}}
\begin{itemize}
  \item Transmitter sends a \textbf{known, fixed-length training sequence} immediately before the
        user data.
  \item Receiver \textbf{knows} that sequence, so it can compute the error exactly.
  \item The adaptive algorithm runs on this sequence and drives the filter taps to their optimum
        values, minimising BER.
  \item Typical sequence: a \textbf{pseudo-random binary signal} or a fixed bit pattern.
  \item Design rule: the training sequence must be \textbf{long enough} for the equalizer to
        converge under the \mk{worst possible} channel conditions.
  \item In GSM the training sequence sits in the \textbf{middle of every burst} (the midamble), so
        the estimate is equally fresh at both ends of the burst. \added
\end{itemize}}{%
\lead{Tracking mode \mk{(follow the channel)}}
\begin{itemize}
  \item Starts the instant training ends and \textbf{user data} begins.
  \item There is no known reference now, so the algorithm uses the \textbf{detected symbols}
        as the reference (decision directed).
  \item It keeps \textbf{tracking} the channel, which is changing as the mobile moves.
  \item \mk{The whole point: a mobile channel changes continuously, so a one-off training would go
        stale within milliseconds.}
\end{itemize}

\lead{Three factors that set the time span an equalizer can handle}
\begin{itemize}
  \item \textbf{Length of the training sequence} and the \textbf{rate of convergence} of the
        algorithm.
  \item \textbf{Speed of the mobile}, which sets the fading rate, the Doppler spread and hence the
        coherence time.
  \item \textbf{Number of taps}, which must cover the maximum expected \textbf{delay spread}. More
        taps $\rightarrow$ more delay spread covered, but \mk{circuit complexity and processing
        time rise with tap count}.
\end{itemize}}

\Q{For the MSE algorithm, determine the optimum solution for the weights}{\m{8}\m{7}}
{\color{sub}\footnotesize \tF{3} \yr{\textbf{\texttt{79 Bh}}, \textbf{79 Ch}, \textbf{80 Ch}}
\gap$\cdot$\gap \mk{asked as a full derivation w/ diagram} $\rightarrow$ worked in the Ch 5
companion}

\sbsr{0.53}{%
\lead{Structure of the adaptive equalizer +Fig}
\begin{itemize}
  \item A \textbf{transversal filter} w/ $N$ delay elements, $N+1$ taps and $N+1$ tunable
        \textbf{complex weights} $w_{0k} \dots w_{Nk}$.
  \item Subscript $k$ $=$ discrete time index, and it appears on the weights to show they are
        \textbf{updated sample by sample}.
  \item Single input $y_k$ at each instant. Its value depends on the instantaneous state of the
        radio channel and on the noise.
  \item Weights are updated continuously by the \textbf{adaptive algorithm}, which is driven by
        the error signal $e_k$.
\end{itemize}}{%
\figT{c5_adaptive_eq.png}
\penalty0\vspace{1pt}
{\color{sub}\footnotesize Basic linear equalizer during training. $d_k$ is the known training
symbol, $e_k$ the error that drives the adaptive algorithm.}}

\sbs{%
\lead{Setting up the cost function}
\begin{itemize}
  \item Input vector and weight vector:
        \[ \mathbf{y}_k = [y_k\;y_{k-1}\;\dots\;y_{k-N}]^T, \qquad
           \boldsymbol{\omega}_k = [w_{0k}\;w_{1k}\;\dots\;w_{Nk}]^T \]
  \item Equalizer output $\hat{d}_k = \mathbf{y}_k^T \boldsymbol{\omega}_k$.
  \item \textbf{Error signal}, w/ $x_k$ the known transmitted symbol:
        \[ e_k = d_k - \hat{d}_k = x_k - \mathbf{y}_k^T \boldsymbol{\omega}_k \]
  \item Squaring:
        \[ e_k^2 = x_k^2 + \boldsymbol{\omega}^T\mathbf{y}_k\mathbf{y}_k^T\boldsymbol{\omega}
                   - 2 x_k \mathbf{y}_k^T \boldsymbol{\omega} \]
  \item \textbf{Mean square error} is the expectation:
        \[ \xi = E[e_k^2]
            = E[x_k^2] + \boldsymbol{\omega}^T \mathbf{R}\, \boldsymbol{\omega}
              - 2\,\mathbf{p}^T \boldsymbol{\omega} \]
        where
        \[ \begin{gathered}
           \mathbf{R} = E\!\left[\mathbf{y}_k \mathbf{y}_k^T\right]
           \text{ (input correlation matrix)}, \\
           \mathbf{p} = E\!\left[x_k \mathbf{y}_k\right]
           \text{ (cross correlation)} \end{gathered} \]
\end{itemize}}{%
\lead{Minimising it}
\begin{itemize}
  \item $\xi$ is a \textbf{quadratic} function of the weights, so its surface is a
        \mk{bowl-shaped paraboloid} (a hyperparaboloid for more than 2 taps).
  \item It is \textbf{concave upward}, so a \textbf{single global minimum} exists.
        \mk{That is why gradient descent cannot get stuck here.}
  \item Take the gradient w.r.t. the weight vector:
        \[ \nabla \xi = \left[\frac{\partial\xi}{\partial w_0}\;
           \frac{\partial\xi}{\partial w_1}\;\dots\;
           \frac{\partial\xi}{\partial w_N}\right]^T
           = 2\mathbf{R}\boldsymbol{\omega} - 2\mathbf{p} \]
  \item Set $\nabla \xi = 0$:
        \[ \boxed{\;\hat{\boldsymbol{\omega}} = \mathbf{R}^{-1}\mathbf{p}\;} \]
        valid as long as $\mathbf{R}$ is \textbf{non-singular}, ie has an inverse.
  \item Substituting back gives the \textbf{minimum MSE}:
        \[ \xi_{\min} = E[x_k^2] - \mathbf{p}^T \hat{\boldsymbol{\omega}}
                      = E[x_k^2] - \mathbf{p}^T \mathbf{R}^{-1}\mathbf{p} \]
  \item \mk{Minimising the MSE tends to minimise the bit error rate}, which is why MSE is the
        standard cost function even though BER is what actually matters.
\end{itemize}

\lead{The iterative form}
\begin{itemize}
  \item Direct matrix inversion is expensive, so in practice the weights are stepped:
        \[ \begin{gathered} \text{new weights} = \text{previous weights} \\
           + \mu \times (\text{previous error}) \times (\text{current input})
           \end{gathered} \]
  \item $\mu$ is a \textbf{step size} adjusted to control how far the weights move per iteration.
  \item Repeated rapidly in a loop until \textbf{convergence}, at which point the weights are
        frozen until the next training sequence.
\end{itemize}}

\hr

\T{5.2 Types of equalizer}

\Q{Explain the structure of the linear transversal equalizer, DFE and MLSE}{\m{6+1+4}\m{4}\m{6+1+5}}
{\color{sub}\footnotesize \tS{6} \yr{71 Ma, 72 Ma, \textbf{72 Ash}, \textbf{\texttt{81 Bh}},
\textbf{\texttt{82 Bh}}} \gap$\cdot$\gap
\mk{MLSE alone is asked in 4 papers, always with a figure}}

\sbsr{0.42}{%
\figT{c5_eq_class.png}
\penalty0\vspace{1pt}
{\color{sub}\footnotesize The full classification: two \textbf{types}, four \textbf{structures},
and the \textbf{algorithms} that drive each one. Reproduce this tree and most of the marks are
already earned.}}{%
\lead{Linear equalizers}
\begin{itemize}
  \item Output is \textbf{not} fed back to change subsequent outputs.
  \item Implemented as an \textbf{FIR filter}: current and past values of the received signal are
        linearly weighted by the tap coefficients and summed.
  \item If the taps are analog, the output is \textbf{sampled at the symbol rate} before the
        decision device.
  \item \textbf{Linear transversal equalizer (LTE)} +Fig: a tapped delay line, delay $T_s$ per
        element, transfer function a function of $z^{-1}$ (or $e^{-j\omega T_s}$).
  \item The simple LTE uses \textbf{feedforward taps only}. A version w/ both feedforward and
        feedback taps exists but is \mk{unstable and rarely used}.
  \item \textbf{Lattice equalizer}: same job, but built as a lattice filter.
  \begin{itemize}
    \item \textbf{($+$)} numerically \textbf{stable}, fast convergence.
    \item \textbf{($+$)} the \mk{order can be changed dynamically} without restarting, so a short
          filter can be used on a benign channel and a long one on a bad channel.
    \item \textbf{($-$)} structurally more complicated.
  \end{itemize}
  \item \mk{($-$) Linear equalizers perform badly on channels with deep spectral nulls}, because
        inverting a null means enormous gain at that frequency, which also amplifies the noise.
        That limitation is the whole reason nonlinear equalizers exist.
\end{itemize}}

\penalty0\vspace{2pt}
\sbs{%
\figT{c5_lte.png}
\penalty0\vspace{1pt}
{\color{sub}\footnotesize Basic linear transversal equalizer: a tapped delay line clocked at the
symbol period $T_s$.}}{%
\figT{c5_ff_fb_taps.png}
\penalty0\vspace{1pt}
{\color{sub}\footnotesize Tapped delay line filter with \emph{both} feedforward and feedback taps.}}

\creamq{Nonlinear equalizers}{\m{4}}
\sbsr{0.55}{%
\begin{itemize}
  \item Used when the channel distortion is \textbf{too severe for a linear equalizer}.
  \item Three effective methods:
        \textbf{DFE}, \textbf{maximum likelihood symbol detection},
        \textbf{MLSE}.
\end{itemize}

\lead{Decision feedback equalizer (DFE)}
\begin{itemize}
  \item \textbf{Basic idea}: once a symbol has been detected and \textbf{decided upon}, the ISI it
        will cause on \textbf{future} symbols can be \textbf{estimated and subtracted out} before
        those symbols are detected.
  \item Two filters:
  \begin{itemize}
    \item \textbf{Feedforward filter (FFF)}: an ordinary linear transversal filter on the received
          signal.
    \item \textbf{Feedback filter (FBF)}: driven by \textbf{previously detected symbols}, and its
          output is subtracted.
  \end{itemize}
  \item Output:
        \[ \hat{d}_k = \sum_{n=-N_1}^{N_2} c_n^{*}\, y_{k-n}
                       + \sum_{i=1}^{N_3} F_i\, d_{k-i} \]
  \item \mk{($+$) The FBF removes ISI without amplifying noise}, because it works on already
        decided symbols, which are noise free. That is exactly what a linear equalizer cannot do.
  \item \textbf{($-$) Error propagation}: one wrong decision feeds a wrong correction into every
        following symbol.
  \item Realised in \textbf{direct transversal} form or as a \textbf{lattice} filter.
  \item \textbf{Predictive DFE} (Belfiore and Park): the FBF acts as a \textbf{noise predictor}
        instead. Performs as well as the conventional DFE in the limit of many taps, and the
        \textbf{RLS} algorithm gives it fast convergence.
\end{itemize}}{%
\figT{c5_dfe.png}
\penalty0\vspace{1pt}
{\color{sub}\footnotesize Decision feedback equalizer: feedforward filter on the received signal,
feedback filter driven by detected symbols.}
\penalty0\vspace{4pt}
\figT{c5_lattice.png}
\penalty0\vspace{1pt}
{\color{sub}\footnotesize Lattice equalizer structure.}}

\creamq{Maximum likelihood sequence estimation (MLSE)}{\m{6}\m{4}}
{\color{sub}\footnotesize \tF{4} \yr{71 Ma, 72 Ma, \textbf{72 Ash}, \textbf{\texttt{81 Bh}}}
\gap$\cdot$\gap \mk{always asked with a block diagram}}

\sbsr{0.55}{%
\begin{itemize}
  \item Ordinary equalizers decode \textbf{each symbol on its own}. MLSE instead
        \textbf{tests all possible data sequences} and picks the sequence with the
        \textbf{maximum likelihood} of having been sent.
  \item \mk{It is optimal in the sense of minimising the probability of a sequence error.}
        No equalizer can do better.
  \item First proposed by \textbf{Forney}, who implemented the MLSE estimator using the
        \textbf{Viterbi algorithm} (see Ch 6).
  \item \textbf{Three blocks} +Fig:
  \begin{itemize}
    \item \textbf{Adaptive matched filter} matched to the combined channel, giving $z(t)$.
    \item \textbf{MLSE} proper, running the Viterbi search over the trellis.
    \item \textbf{Channel estimator} in a feedback loop, keeping the filter and the metric
          matched to the current channel.
  \end{itemize}
  \item \textbf{Requirements:}
  \begin{itemize}
    \item knowledge of the \textbf{channel characteristics}, to compute the branch metrics,
    \item knowledge of the \textbf{statistical distribution of the noise}.
  \end{itemize}
  \item \mk{($-$) Huge computational load}: complexity grows \textbf{exponentially} w/ the channel
        delay spread in symbols, so it is only practical when the delay spread covers a small
        number of symbols. GSM is exactly that case.
\end{itemize}}{%
\figT{c5_mlse.png}
\penalty0\vspace{1pt}
{\color{sub}\footnotesize MLSE with an adaptive matched filter and a channel estimator in the
loop.}}

\Q{Explain any two adaptive equalization algorithms}{\m{4}\m{5}}
{\color{sub}\footnotesize \yr{73 Ma, 74 Ma} \gap$\cdot$\gap \mk{name all three, then detail two}}

\penalty0\vspace{2pt}
\begin{tabularx}{\textwidth}{@{}L{2.6cm}XXX@{}}
\toprule
\hrow & \thead{Zero forcing (ZF)} & \thead{Least mean square (LMS)} & \thead{Recursive least
squares (RLS)} \\
\midrule
Criterion & Force ISI to \textbf{zero} at the sampling instants & Minimise the \textbf{mean square
error} & Minimise a \textbf{weighted sum of past squared errors} \\
Update & Taps set so the combined response is 1 at $t=0$ and 0 at every other $kT_s$ &
  $\boldsymbol{\omega}_{k+1} = \boldsymbol{\omega}_k + \mu\, e_k\, \mathbf{y}_k^{*}$ &
  Uses the Kalman gain vector and an inverse correlation matrix \\
Complexity & Lowest & Low, \textbf{$2N+1$} operations per iteration & \textbf{High},
  $\propto N^2$ \\
Convergence & Immediate (non-adaptive form) & \mk{Slow}, needs a long training sequence &
  \mk{Fastest}, an order of magnitude fewer iterations \\
\midrule
What sets the speed \added & Fixed by the channel & \mk{The eigenvalue spread}
  $\lambda_{max}/\lambda_{min}$ of $\mathbf{R}$: a large spread $\rightarrow$ the error bowl is a
  long narrow valley $\rightarrow$ one step size $\mu$ cannot suit both axes &
  \mk{Insensitive to the eigenvalue spread}, which is exactly why it wins on a
  badly conditioned channel \\
Noise & \mk{Amplifies noise badly} at spectral nulls & Well behaved & Well behaved \\
Use it when & The channel is benign and simple hardware matters & Convergence speed is not
  critical & The channel changes fast, or training must be short \\
\bottomrule
\end{tabularx}

\penalty0\vspace{3pt}
\begin{itemize}
  \item \textbf{Zero forcing in more detail:} it is the \textbf{inverse filter} of the channel,
        $H_{eq}(f) = 1/H_{ch}(f)$. It removes ISI completely \mk{in a noise-free channel}.
  \item \mk{($-$) In reality it does not perform well}: where the channel has a spectral null,
        $1/H_{ch}(f)$ becomes huge, so the equalizer applies enormous gain exactly where the noise
        is unopposed. It \textbf{amplifies noise} and can perform worse than no equalizer at all.
  \item \textbf{How to choose:} the \textbf{speed of the mobile} sets the fading rate and Doppler
        spread, hence the coherence time. The algorithm must converge well inside that coherence
        time, so a fast-moving mobile forces RLS despite its cost.
\end{itemize}

\hr

\T{5.3 Diversity}

\Q{Why is diversity needed? Explain the different types of diversity technique}{\m{2+6}\m{4+6}\m{2+2}\m{3+3+2}\m{2}}
{\color{sub}\footnotesize \tS{11} \yr{\textbf{70 Bh}, \textbf{71 Bh}, 71 Ma, 73 Ma,
\textbf{74 Bh}, 74 Ma, \textbf{76 Bh}, \textbf{78 Ch}, \texttt{80 Ba}, \textbf{\texttt{80 Bh}},
\texttt{81 Ba}} \gap$\cdot$\gap \mk{second most repeated question in the chapter}}

\sbsr{0.62}{%
\lead{The concept \mk{(the ``why'' marks live here)}}
\begin{itemize}
  \item \textbf{Objective}: combine multiple copies of the signal so as to
        \mk{reduce the effect of deep fades}.
  \item \textbf{The core idea in one sentence}: if one radio path is in a deep fade, an
        \textbf{independent} path probably is not. Having more than one path to choose from
        improves both the instantaneous and the average SNR.
  \item Typical improvement: \mk{20 to 30 dB}.
  \item \textbf{($+$) Requires no training overhead}, unlike equalization.
  \item \textbf{($+$)} Large link improvement for little added cost.
  \item Diversity decisions are made \textbf{entirely at the receiver} and are unknown to the
        transmitter.
  \item \mk{Equalization fixes ISI. Diversity fixes fading. They solve different problems and are
        often used together.}
\end{itemize}}{%
\figT{c5_snr_div.png}
\penalty0\vspace{1pt}
{\color{sub}\footnotesize Probability distribution of SNR threshold for $M$-branch selection
diversity. $\Gamma$ is the mean SNR on each branch. \mk{Each extra branch pushes the low-SNR
tail sharply to the right, which is the diversity gain.}}}

\sbs{%
\lead{Microscopic Vs macroscopic diversity}
\begin{itemize}
  \item \textbf{Microscopic diversity} counters \textbf{small-scale} fading.
  \begin{itemize}
    \item Selects the best signal at all times out of several rapidly fading ones.
    \item Antennas separated by a fraction of a wavelength are enough.
  \end{itemize}
  \item \textbf{Macroscopic diversity} counters \textbf{large-scale} fading, ie shadowing.
  \begin{itemize}
    \item Uses \textbf{base stations that are not shadowed}, sufficiently separated in space.
    \item Eg CDMA soft handoff, where two BSs serve one mobile.
  \end{itemize}
\end{itemize}}{%
\lead{Types of diversity \mk{(name all five)}}
\begin{itemize}
  \item \textbf{Space (antenna) diversity}
  \begin{itemize}
    \item Two or more \textbf{physically separated} receiving antennas.
    \item Separation of \textbf{$\lambda/2$ or more} at the mobile gives essentially
          \textbf{uncorrelated envelopes}.
    \item At a BS the separation must be much larger, \mk{up to 20 wavelengths}, because the
          BS antenna is high and sees a narrow angular spread.
    \item \mk{The most common form, and the one the examiner means when the type is unspecified.}
  \end{itemize}
  \item \textbf{Polarization diversity}
  \begin{itemize}
    \item Send or receive the same information on \textbf{two orthogonal polarizations}, normally
          vertical and horizontal.
    \item Relies on the \textbf{decorrelation of the two receive ports}, so the two ports must
          stay cross-polarized.
    \item \mk{($+$) Needs no antenna spacing at all}, so one dual-polarized antenna replaces two
          widely spaced ones. Huge saving on BS tower space.
    \item \textbf{($-$)} Only 2 branches are possible, and each port gets about
          \textbf{3 dB less power} than a single co-polarized antenna.
  \end{itemize}
\end{itemize}}

\sbs{%
\begin{itemize}
  \item \textbf{Frequency diversity}
  \begin{itemize}
    \item Transmit the same information on \textbf{more than one carrier frequency}.
    \item Carriers separated by \textbf{more than the coherence bandwidth} do not fade together.
    \item \textbf{($-$)} Costs extra spectrum and an extra transmitter chain.
    \item Used on microwave LOS links, where it doubles as equipment redundancy.
  \end{itemize}
\end{itemize}}{%
\begin{itemize}
  \item \textbf{Time diversity}
  \begin{itemize}
    \item Repeatedly transmit the information at time spacings \textbf{exceeding the coherence
          time} of the channel.
    \item Repetitions then fade independently.
    \item \textbf{Two implementations \mk{(72 Ash asks for exactly two, \m{5})}}:
    \begin{itemize}
      \item \textbf{Interleaving with channel coding}: spread a codeword's bits far apart in time
            so a burst error becomes several correctable single errors.
      \item \textbf{RAKE receiver} in CDMA: multipath copies are themselves time-shifted versions
            of the same signal, so combining them \mk{is time diversity for free}.
      \item \textbf{ARQ} retransmission is a third, cruder form. \added
    \end{itemize}
    \item \textbf{($-$)} Costs delay, which \mk{voice cannot tolerate beyond about 40 ms}.
  \end{itemize}
\end{itemize}}

\Q{Explain the combining methods used for space diversity}{\m{4+4}\m{3+3+2}\m{2+2}\m{4}}
{\color{sub}\footnotesize \tS{7} \yr{74 Ma, 73 Ma, \textbf{74 Bh}, 71 Ma, \textbf{78 Ch},
\texttt{80 Ba}, \textbf{\texttt{80 Bh}}} \gap$\cdot$\gap 78 Bh asks specifically for
\mk{feedback Vs maximal ratio combining} \gap$\cdot$\gap 80 Bh asks for MRC alone}

\sbsr{0.42}{%
\figT{c5_space_div.png}
\penalty0\vspace{1pt}
{\color{sub}\footnotesize Generalized block diagram for space diversity: $m$ antennas, variable
gains, then switching logic or a demodulator.}}{%
{\color{sub}\footnotesize Space diversity reception falls into \textbf{four} categories. They
differ only in \textbf{how the $m$ branches are weighted and combined}, so answer them as one
family. The general model has $m$ branches whose gains are adjusted to give the same average SNR
per branch.}
\penalty0\vspace{2pt}
\begin{tabularx}{\linewidth}{@{}L{2.35cm}XL{1.85cm}@{}}
\toprule
\hrow \thead{Method} & \thead{How it works} & \thead{Cost} \\
\midrule
\textbf{Selection} & Monitor SNR on all $m$ branches, connect \textbf{the highest} to the
  demodulator & $m$ receivers, needs all channels monitored \\
\textbf{Feedback / scanning} & Scan branches in a \textbf{fixed sequence}, take the first above a
  \textbf{preset threshold}, keep it until it falls below & \mk{Only one receiver} \\
\textbf{Maximal ratio (MRC)} & \textbf{Co-phase} all branches, weight each by its own
  \textbf{signal-voltage to noise-power ratio}, then sum & $m$ receivers, needs accurate gain
  estimates \\
\textbf{Equal gain (EGC)} & \textbf{Co-phase} all branches, weight every one by \textbf{unity},
  then sum & $m$ receivers, no gain estimate \\
\bottomrule
\end{tabularx}}

\penalty0\vspace{3pt}
\sbsr{0.55}{%
\lead{Selection diversity}
\begin{itemize}
  \item \textbf{Principle:} pick the best signal at every instant.
  \item Simplest in concept, but \mk{requires knowledge of all channels} $\rightarrow$ $m$
        monitoring receivers.
  \item \textbf{Performance:} for $M$ branches w/ mean SNR $\Gamma$ each,
        \[ \Pr[\text{SNR} \le \gamma] = \left(1 - e^{-\gamma/\Gamma}\right)^{M} \]
        and the average SNR gain is
        \[ \frac{\bar{\gamma}_M}{\Gamma} = \sum_{k=1}^{M}\frac{1}{k} \]
  \item \mk{Diminishing returns}: going from 1 to 2 branches gains 1.5$\times$, 2 to 3 only
        1.83$\times$. \textbf{Most of the benefit is in the second branch.}
\end{itemize}

\lead{Feedback or scanning diversity}
\begin{itemize}
  \item \textbf{Principle:} the $M$ signals are scanned in a \textbf{fixed sequence} until one is
        found \textbf{above a predetermined threshold}.
  \item That signal is used \textbf{until it falls below threshold}, then scanning restarts.
  \item \textbf{($+$) Simple to implement, only one receiver is required.} That is its whole
        selling point.
  \item \textbf{($-$)} The resulting fading statistics are \mk{somewhat inferior} to the other
        three, because the branch in use is merely acceptable, not best.
\end{itemize}}{%
\figT{c5_selection.png}
\penalty0\vspace{2pt}
\figT{c5_scanning.png}
\penalty0\vspace{1pt}
{\color{sub}\footnotesize Top: selection diversity, monitor SNR then select highest. Bottom:
scanning diversity, one receiver plus a comparator against a preset threshold.}}

\penalty0\vspace{3pt}
\sbsr{0.55}{%
\lead{Maximal ratio combining (MRC) \mk{(80 Bh asks for this alone, \m{2})}}
\begin{itemize}
  \item \textbf{Principle:} combine \textbf{all} the signals, \textbf{co-phased} and
        \textbf{weighted}, so as to have the highest achievable SNR at all times.
  \item Branches are weighted by their \textbf{signal voltage to noise power ratio}, then summed.
  \item Co-phasing first is essential, otherwise the branches would partly cancel.
  \item \textbf{Output SNR is the sum of the individual SNRs:}
        \[ \gamma_M = \sum_{i=1}^{M}\gamma_i \]
  \item \mk{So MRC can produce an acceptable output even when none of the individual branches is
        acceptable on its own.} That is the line the examiner wants.
  \item \textbf{($+$) Gives the best statistical reduction of fading} of any linear combiner.
  \item \textbf{($-$)} Needs an \textbf{accurate estimate of each branch's amplitude gain}, so the
        receiver is the most complex of the four.
\end{itemize}}{%
\figT{c5_mrc.png}
\penalty0\vspace{1pt}
{\color{sub}\footnotesize Maximal ratio combiner: weight, co-phase, sum, then detect.}}

\sbsr{0.55}{%
\lead{Equal gain combining (EGC)}
\begin{itemize}
  \item Used when it is \textbf{not convenient} to provide the variable weighting MRC needs.
  \item Branch weights are \textbf{all set to unity}, but the signals are still \textbf{co-phased}.
  \item Lets the receiver exploit signals received simultaneously on every branch.
  \item \mk{Performance is only marginally inferior to MRC, and superior to selection diversity},
        for far less hardware. That trade is why EGC is common in practice.
\end{itemize}}{%
\figT{c5_egc.png}
\penalty0\vspace{1pt}
{\color{sub}\footnotesize Equal gain combining: unity weights, co-phasing and summing only.}}

\figC{c5_polarization.png}{0.46\textwidth}
{\color{sub}\footnotesize Theoretical model for base station polarization diversity: two
orthogonal receive ports at $\pm\alpha$ in the $x$--$y$ plane.}

\hr

\T{5.4 RAKE receiver and interleaving}

\Q{Explain the working of an M-branch RAKE receiver with a block diagram}{\m{8}\m{3+5}\m{4}\m{3}}
{\color{sub}\footnotesize \tS{6} \yr{70 Ma, \textbf{72 Ash}, \textbf{73 Bh}, \textbf{76 Bh},
\textbf{77 Ch}, \texttt{79 Bh}} \gap$\cdot$\gap 73 Bh asks
\mk{how it exploits time diversity}, so say that explicitly}

\sbsr{0.53}{%
\lead{What it is}
\begin{itemize}
  \item A \textbf{diversity receiver designed specifically for CDMA}.
  \item In an ordinary receiver multipath is a problem. \mk{In a RAKE receiver multipath is the
        resource.}
  \item It \textbf{collects the time-shifted versions} of the original signal by providing a
        \textbf{separate correlation receiver} for each multipath component.
  \item Each correlator is called a \textbf{finger}. Hence ``rake'': fingers gathering energy
        that would otherwise be wasted.
\end{itemize}

\lead{Why it works only in CDMA \mk{(the key insight)}}
\begin{itemize}
  \item The multipath components are \textbf{practically uncorrelated} from one another once
        their relative propagation delays exceed \textbf{one chip period}.
  \item That comes from the \textbf{low autocorrelation} of the PN spreading sequence: a copy
        delayed by more than a chip correlates almost to zero.
  \item So a correlator locked to delay $\tau_m$ sees component $m$ clearly and treats the others
        as noise.
  \item \mk{A narrowband system cannot do this} because its copies are not separable in delay.
\end{itemize}

\lead{How it exploits time diversity \mk{(73 Bh, \m{3})}}
\begin{itemize}
  \item Time diversity means receiving the same information at times separated by more than the
        channel's coherence time or, here, resolvable in delay.
  \item The multipath copies \textbf{are} exactly that: the same symbol, arriving at different
        delays, faded \textbf{independently}.
  \item So the RAKE gets time diversity \mk{without spending any extra bandwidth, power or time},
        which no other diversity scheme manages.
\end{itemize}}{%
\figT{c5_rake.png}
\penalty0\vspace{1pt}
{\color{sub}\footnotesize An $M$-branch ($M$-finger) RAKE. Each correlator detects one time-shifted
version of the CDMA transmission, delayed by at least one chip from the others.}

\lead{Operation, step by step}
\begin{itemize}
  \item \textbf{Multiple correlators} separately detect the $M$ \textbf{strongest} multipath
        components. Correlator $m$ is delay-locked to component $m$ and outputs $Z_m$.
  \item Outputs are \textbf{weighted} by $\alpha_m$ to produce a better estimate than any single
        component could.
  \item \textbf{Combine} (this is MRC applied to multipath):
        \[ Z' = \sum_{m=1}^{M}\alpha_m Z_m, \qquad
           \alpha_m = \frac{Z_m^{2}}{\displaystyle\sum_{m=1}^{M} Z_m^{2}} \]
  \item \textbf{Demodulation and bit decision} are then made on the weighted sum $Z'$, not on any
        one finger.
  \item If one finger is in a deep fade, the others carry the symbol.
  \item \mk{A CDMA link in a rich multipath environment can therefore outperform the same link in
        a ``clean'' one-path environment.}
\end{itemize}}

\Q{What is interleaving? Explain its working principle}{\m{4}\m{2+6}}
{\color{sub}\footnotesize \tF{3} \yr{\textbf{\texttt{80 Bh}}, \textbf{\texttt{81 Bh}},
\texttt{81 Ba}}}

\sbsr{0.55}{%
\lead{Why it is needed}
\begin{itemize}
  \item A fade does not corrupt bits at random. It wipes out a \textbf{burst} of consecutive bits.
  \item Channel codes (Ch 6) are designed to fix \textbf{scattered} errors. A burst overwhelms
        them.
  \item \textbf{Interleaving scrambles the time order of the source bits before channel coding},
        so that a burst on the channel arrives at the decoder as \mk{isolated single errors spread
        across many codewords}, which the code can fix.
  \item Speech coders make this worse: several \textbf{``important'' bits} occur in succession, so
        without interleaving one fade can destroy an entire speech frame.
  \item \mk{Interleaving is time diversity implemented in the digital domain.}
\end{itemize}}{%
\figT{c5_interleaver.png}
\penalty0\vspace{1pt}
{\color{sub}\footnotesize Block interleaver: coded bits are read \textbf{in} by column and read
\textbf{out} to the modulator one row at a time.}}

\sbsr{0.55}{%
\lead{Block interleaver \mk{(the figure above)}}
\begin{itemize}
  \item Formats the encoded data into a rectangular array of \textbf{$m$ rows} and
        \textbf{$n$ columns}, interleaving $nm$ bits at a time.
  \item Each row usually holds one \textbf{$n$-bit source word}.
  \item An interleaver of \textbf{degree (depth) $m$} has $m$ rows.
  \item \textbf{Write in by column}: bits are placed by increasing the row number for each
        successive bit, filling the columns.
  \item \textbf{Read out by row}, and transmit.
  \item Effect: \mk{original neighbouring bits end up $m$ bit periods apart}.
  \item At the receiver the \textbf{de-interleaver} does the inverse: store by column, clock out
        one row at a time.
  \item So an $(n,k)$ code that can handle a burst of $b < (n-k)/2$ can be combined w/ an
        interleaver of degree $m$ to make an $(mn, mk)$ block code that handles bursts of
        \textbf{$mb$} bits.
\end{itemize}}{%
\lead{Convolutional interleaver}
\begin{itemize}
  \item Same job, but built as a bank of shift registers of increasing length rather than a
        rectangular array.
  \item Can replace a block interleaver in the same fashion, and is \textbf{ideally suited for use
        with convolutional codes}.
  \item \textbf{($+$)} Roughly \mk{half the delay and half the memory} of a block interleaver of
        the same effective depth. \added
\end{itemize}}

\lead{The cost}
\begin{itemize}
  \item \textbf{Delay}, and the delay grows with interleaver depth.
  \item \mk{Human speech tolerates delays below about 40 ms}, which caps the depth a voice system
        can use. Data services can interleave far more deeply.
\end{itemize}
